\begin{abstract}
Large language models are increasingly deployed in systems that require conditional behavior---customer service workflows, content moderation pipelines, and agentic tool-use chains---where system-prompt rules take the form ``if condition $X$ holds, then do $Y$.'' When these conditional instructions misfire, the consequences range from degraded user experience to safety violations. Yet no existing benchmark systematically isolates and measures in-context if-then capacity across different condition types. We introduce \benchname, a programmatic benchmark of 144 test cases spanning 7 condition categories---lexical, semantic, near-miss, sustained/adversarial, scaling, negation, and if-then-else branching---with fully deterministic verification via distinctive marker strings. We evaluate three models from the GPT-4.1 family (large, mini, nano) and find that accuracy varies significantly by condition type ($\chi^2$ test, $p < 0.002$ for all models), with sustained/adversarial conditions proving hardest (42.9--78.6\%) and if-then-else branching easiest (90--100\%). Across all models and conditions, false positives outnumber false negatives by 2.8:1 ($p < 0.001$), revealing a systematic bias toward over-triggering conditional rules. Model size strongly predicts performance: \gptlarge achieves 97.2\% case-level accuracy versus 87.5\% for \gptmini and 68.1\% for \gptnano. These results provide the first fine-grained characterization of where and how LLMs fail at conditional instruction following.
\end{abstract}
